\scanpage{360}\mspage{276}
On suppose que $N$ satisfait $L_{\pi_0}$ (stabilité
par $\unsure{\underline{\lim}|\underrightarrow{\lim}}$ gdes devant $\pi_0$). Soit\deleted{\unsure{nt}} $M=\mathrm{Ind}_{\pi_0}C$,
$\underline{\mathrm{Hom}}_{!(\pi_0)}(M,N)$ la sous-catégorie \inserted{str.} pleine
de $\underline{\mathrm{Hom}}(M,N)$ formée des $f:M\unsure{-|\longrightarrow}N$
\deleted{qui \unsure{commutent aux}}\inserted{$\pi_0$-accessibles}. Alors le foncteur
restriction
\[
 \underline{\mathrm{Hom}}_{!\deleted{\unsure{(}}\pi_0\deleted{\unsure{)}}}(M,N)
 \longrightarrow\underline{\mathrm{Hom}}(C,N) \tag*{(4.12.3.1)}
\]
est une équivalence de catégories.
On \unsure{trouve|donne} un foncteur quasi-inverse
comme composé
\[
\begin{aligned}
 \underline{\mathrm{Hom}}(C,N)&\longrightarrow
 (\underline{\mathrm{Hom}}\,\underbrace{\mathrm{Ind}_{\pi_0}(C)}_{M},\mathrm{Ind}_{\pi_0}(N))\\
 &\longrightarrow\underline{\mathrm{Hom}}(M,N),
\end{aligned}\tag*{(4.12.3.2)}
\]
où la deuxième flèche est le foncteur
composition $F\to L\circ F$ avec le foncteur
limite inductive $L:\mathrm{Ind}_{\pi_0}N\unsure{-|\longrightarrow}N$.
Le foncteur composé (4.12.3.2) \deleted{\unsure{qui aura}}
a une image contenue dans $\underline{\mathrm{Hom}}_{!\pi_0}(M,N)$
et induit une équivalence entre \inserted{cette dernière catégorie et}
$\underline{\mathrm{Hom}}(C,N)$. \deleted{\unsure{et celle}}
% archive sheet number: 359

